\section{Zadanie 11}
Losujemy słowo binarne długości \(n\). Rozważmy wszystkie \(\binom{n}2\)
par indeksów. Dla każdej takiej pary zmienna \(X_{ij}\) jest wartością
funkcji xor na \(i\)-tym oraz \(j\)-tym bicie. Niech X = 
\(\sum_{1\leqslant i < j \leqslant n} X_{ij}\).

\begin{enumerate}[label=(\roman*)]
  \item wykaż, że zmienne nie są niezależne jako rodzina zmiennych
  \item wykaż, że \(\mathbb{E}(X_{ab}X_{cd}) = 
    \mathbb{E}(X_{ab}) \mathbb{E}(X_{cd})\) dla 
    \(\{a, b\} \neq \{c, d\}\)
  \item wyznacz Var(\(X\)) korzystając z poprzedniego podpunktu
  \item przy pomocy nierówności Czebyszewa ogranicz
  \(P(|X-E(X)|\geqslant n)\)
\end{enumerate}

\subsection*{Rozwiązanie}
\subsubsection*{(i)}
Zakładamy \(n \geqslant 3\) (dla mniejszego \(n\) jest \(\leqslant 1\) 
zmiennych losowych, więc na pewno są niezależne)

Oczywiście \(P(X_{ij} = 0) = P(X_{ij} = 1) = \frac{1}{2}\quad
\forall_{1\leqslant i < j \leqslant n}\)

Weźmy trzy bity \(A_1, A_2\) i \(A_3\). 

Zdarzenie \((X_{12} = 1 \land X_{13} = 1 \land X_{23} = 1)\)
wymaga \(A_1 \neq A_2, A_1 \neq A_3\) i \( A_2 \neq A_3\),
co nie jest możliwe jednocześnie.

Otrzymujemy \(P(X_{12} = 1 \land X_{13} = 1 \land X_{23} = 1) = 0
\neq \frac{1}{8} = P(X_{12} = 1) P(X_{13} = 1) P(X_{23} = 1)\) 

\subsubsection*{(ii)}
Jeśli \(\{a, b\} \cap \{c, d\} = \emptyset\) to trywialne - \(X_{ab}\) i \(X_{cd}\)
są niezależne.

Wpp. bez straty ogólności niech \(a = c\), \(b \neq d\).
Wartości bitów \(A_a, A_b, A_d\) są niezależne i z równym prawdopodobieństwem
równe \(0\) lub \(1\).

Oczywiście
\(E(X_{ab}) = P(X_{ab} = 1) = P(A_a \neq A_b) = \frac{1}{2}
= E(X_{ad}) = E(X_{bd})\)

\(E(X_{ab}X_{ad}) = P(X_{ab} = 1 \land X_{ad} = 1) = P(A_a \neq A_b \land 
A_a \neq A_d)\)\\
\(= P(A_b = A_d = 1 | A_a = 0)P(A_a = 0) + P(A_b = A_d = 0 | A_a = 1)P(A_a = 1)
\)

Skoro \(A_a, A_b, A_d\) są niezależne, to można podstawić wszędzie
\(\frac{1}{2}\)

\(E(X_{ab}X_{ad}) = \left(\frac{1}{2}\right)^3
+ \left(\frac{1}{2}\right)^3 = \frac{1}{4}
= E(X_{ab})E(X_{ad})
\)

\subsubsection*{(iii)}
Skoro każda para zmiennych jest niezależna, to wariancja sumy jest
sumą wariancji.
\[\text{Var}(\sum_{1\leqslant i < j \leqslant n} X_{ij})
= \sum_{1\leqslant i < j \leqslant n} \text{Var}(X_{ij})
= \binom{n}{2} \cdot \frac{1}{4}
\]

\subsubsection*{(iv)}
\[
    P(|X-E(X)|\geqslant n)\) \leqslant \frac{\text{Var}(X)}{n^2}
    = \frac{\frac{1}{8}n(n-1)}{n^2} = \frac{n-1}{8n} < \frac{1}{8}
\]

\subsection*{Odpowiedź (iv)}
\[
\boxed{\frac{n-1}{8n}}
\]
